\section{Conclusion}
We evaluated \nEncoders{} frozen vision foundation backbones at \nEncoderConfigs{}
encoder-resolution configurations, under \nHeads{} readouts, across \nDatasets{} person and
vehicle re-identification datasets and \nProtocols{} exact protocols, for \nRuns{} run records
in total, without training a single backbone.

The result we would most like to see acted on is methodological. \textbf{The readout is an
experimental axis, and it is a larger one than the backbone differences that frozen-encoder
comparisons exist to detect.} On Market-1501, changing only how cached features are read moves
SigLIP2-g from rank \rankSiglipgFrozen{} to rank \rankSiglipgProbed{}, and DINOv3-H+ from rank
\rankDinohpFrozen{} to rank \rankDinohpProbed{}. Contrastively pre-trained encoders gain least from a fitted head because
their objective already optimised the geometry that zero-shot cosine retrieval reads, so
frozen-cosine benchmarking flatters them systematically. A matched label-free control shows the
gain is the supervision and not the 512-dimensional bottleneck. Any comparison of frozen
encoders that does not name its readout --- along with its resize convention and input size,
each of which we show to be worth more than several of the differences being reported --- is
under-specified, and can report the reverse of the truth with a table that looks correct.

On the specific question of agglomerative distillation, the answer is two-sided and neither
side is the one the concern anticipated. Where the probe has labels, distillation does not cost
the instance margin --- it improves it, by a wider factor on mINP (\minpGap{}$\times$) than on
mAP. Where the probe does not have labels, the student regresses toward the mean of teachers
that disagree sharply by domain, and the size of that regression tracks the size of the
disagreement: two independent datasets at a teacher spread near 3$\times$ both cost the
student close to 45\% relative. That makes the cost of agglomeration predictable
from a question a practitioner can answer without running the experiment. The practical
recommendation follows and needs its qualifier attached: an agglomerative backbone is the right
default \emph{if you can fit a head on labelled data from your own domain}; if you cannot, the
better single teacher for your domain is available and is better.

Finally, a frozen agglomerative backbone with one 512-dimensional affine head reaches
\bestMarketMAP{} mAP and \bestMarketRone{} Rank-1 on \texttt{market1501/official@1}, for one
forward pass over cached crops plus well under a minute of fitting on a consumer GPU. It does
not beat a tuned specialist. It is a reminder of how much of a modern ReID number is already
present in a general-purpose representation, waiting on the readout.

The natural completion of this work is \cref{sec:decisions}: to ask the same encoders and heads
not how well they rank, but how well they decide --- whether a match should be accepted at all,
at a threshold fixed in advance, against a gallery that may not contain the answer. The
substrate that produced every table here already expresses those measures.

\section*{Reproducibility}
The evaluation package, the protocol definitions with their digests, the encoder and head
specifications, and the complete run records --- including per-query results, the provenance of
every row, and the licence index --- are released under the MIT licence. The tables in this
paper are generated from those records by a script included with the source; no metric in this
document was transcribed by hand. We redistribute no dataset; obtaining each one, and complying
with its licence, remains the reader's responsibility.
